\section{Regular lens ends and the final boundary count}
\label{sec:gluing}

Fix the finite stack, its determinant transitions, its metrics, and the
small perturbations obtained from Propositions~\ref{prop:analysis:first-stage}
and \ref{prop:analysis:finite-induction}.  These choices exclude the
forbidden fixed and mixed candidates and make the trace-main problems
regular.  We construct their lens collars before completing the
regularization of the unbroken free stratum.

All artificial isolations and all surgery seams have finite length.
Only the primary three-sphere cuts and the lens cuts are boundary
faces of the parameter space.  The only end types left by the exclusions
are the legitimate unbroken instantons and a single separated minimal
trace cap with a free main component.  The latter limits have no
particles, cylinder losses, additional true cuts, or corners.
The instanton collars are supplied by
Proposition~\ref{prop:analysis:instanton-link}; we now construct and
exhaust the remaining lens collars.

\subsection{The framed minimal cap}

Let $N$ be the cylindrical completion of the disk bundle of Euler number
$-4$, with zero section $S$, and fix a representative $\gamma$ of the
trace-zero limiting flat connection.  Its determinant-one parallel group
is
\[
 H=\{\diag(e^{i\theta},e^{-i\theta}):\theta\in\R\}
 \cong\U(1).
\]
The cap $N$ has exactly one matching end.  On $N$, \emph{framed}
means quotienting by gauge transformations tending
to the identity at this fixed representative.  Thus $H$ still acts on the
framed configuration space.  On the main component we instead allow
parallel limiting gauge transformations.  The determinant connections
and all identifications in this distinction are fixed.

\begin{proposition}[Regular minimal cap]
\label{prop:gluing:regular-cap}
At each of finitely many regular isolated main-component limits on a
trace face, the cap data can be chosen arbitrarily small, compatibly with
the preceding exclusions, so that the following statements hold.
The framed cap problem of charge $1/4$, with its sphere-sweep mark and
the normalized holonomy condition $-1$, is a finite set of regular
points.  Every point is a crossing on the unique harmonic Abelian
connection with ordered summand difference $v(S)=2$.  These points are
fixed by $H$, and their oriented sum has absolute value one.
The spinor is zero, and the coupled Dirac operator is invertible at
each point.  The choices for the two determinant lifts are related by
the prescribed cap tensor operation.
\end{proposition}

\begin{proof}
The boundary eigenline of the ordered splitting is prescribed by its
degree modulo four.  With that eigenline fixed, the Abelian ASD extension
is unique: the harmonic curvature is determined by the relative class,
the difference between two such connections is a closed one-form modulo
gauge, and $H^1(N;\R)=0$.  The vanishing of $b^+(N)$ gives surjectivity
of its tangential ASD deformation problem.  A different parallel
boundary identification is absorbed by the parallel stabilizer of this
same connection.  In particular, the two eigenlines are not an extra
pair of framed cap solutions.

Corollary~\ref{cor:indices:trace-minimum} gives the following dimensions
before the sweep cut: the tangential framed ASD part has real index
zero, the normal ASD part has complex index one, and the coupled Dirac
part has complex index zero.  Positive scalar curvature and the
specified ASD fixed-line curvature kill the Dirac kernel; its index
then kills its cokernel.  The same Weitzenb\"ock argument kills every
cap spinor for the allowed small curvature perturbations.  There is no
Dirac forcing on the separated cap.

We explain why the two remaining normal regularity requirements can
be imposed by the allowed perturbations.  Write $T$ for the complex
normal ASD operator at the reducer.  Its kernel and cokernel are
finite-dimensional complex representations of $H$, all with the
off-diagonal character.  Derivatives of based holonomies detect
$\ker T$ modulo infinitesimal gauge.  Indeed, suppose a deformation has
on every based loop only the derivative produced by the same
infinitesimal conjugation.  Subtract that conjugation and compare
parallel transport from the base point along paths.  Independence of
the path, which follows from the vanishing loop derivatives, defines
an infinitesimal gauge transformation producing the deformation.
Its off-diagonal part tends to zero on the end, since $\gamma$ has
no off-diagonal parallel section.  It is therefore a framed gauge
direction.  Consequently the holonomy derivatives separate the normal
kernel.  Finitely many suffice: successively reduce their common
kernel, whose dimension is finite.

Include a noncentral holonomy in this finite sample to fix the torus.
In a slice for the sample orbit, a torus-equivariant linear function
of the normal sample coordinates, followed by a localized
off-diagonal two-form output, vanishes on commuting data and has the
prescribed normal derivative.  Multiplication by an invariant cutoff
extends it to an allowed perturbation.  Smooth compactly supported
outputs detect every nonzero cokernel element, so finitely many
outputs suffice as well.  These constructions realize all complex
linear maps from $\ker T$ to $\coker T$.

For completeness, split domain and range into kernel, cokernel, and
complements on which $T$ is invertible.  The Schur complement of a
small perturbation is a finite-dimensional map
$\ker T\longrightarrow\coker T$.  Surjective such maps are dense
because $\dim_{\C}\ker T-\dim_{\C}\coker T=1$.
The preceding perturbations supply all their first derivatives.
We may thus make $T$ surjective by an arbitrarily small perturbation
whose value on the reducer remains zero.  Its kernel is now a complex
line.

In the determinant-root normalization the reducer's sphere sweep is
a loop in $H$.  The eigenline has degree $v(S)/2=1$, so the sweep has
winding $1$ or $-1$, according to its fixed orientation.  Perturb its
toral component slightly, relative to the collapsed ends, so that its
crossings of $-1$ are transverse.  Their signed sum is this winding.
At a crossing, left translation identifies the target tangent space
with
\[
 T_{-1}\SU(2)=\R\tau\oplus\C_{\mathrm{off}},
\]
where $\tau$ generates the diagonal Lie algebra.  The sweep-mark
derivative is nonzero in $\R\tau$.  On the complex normal kernel, the
off-diagonal cut derivative is $H$-equivariant and hence complex linear.
The same finite sample construction, now with output in the cut
target, varies this map arbitrarily.  Choose it nonzero at every
crossing.  Equation and cut variations are independent, so this does
not undo surjectivity of $T$.  The combined cut derivative on the
kernel and mark is therefore an isomorphism
\begin{equation}
 \ker T\oplus\R\ \longrightarrow\
 \C_{\mathrm{off}}\oplus\R\tau.
 \label{eq:gluing:cap-cut}
\end{equation}
Its complex factor preserves orientation, leaving precisely the
toral crossing sign.  The framed dimension calculation here is
\begin{equation}
 2+1-3=0.
 \label{eq:gluing:cap-dimension}
\end{equation}
This is different from the ordinary interval calculation
$1+1+1-3=0$: the latter uses the unframed index of $W'$, an interval
parameter, and a sweep mark.  No interval parameter is being added
to the fixed framed cap in \eqref{eq:gluing:cap-dimension}.

It remains to exclude additional cap points.  Near every reducer
crossing, \eqref{eq:gluing:cap-cut} and the implicit function theorem
give a unique framed solution for the fixed tangential data.  It is
the reducer crossing just constructed.  Away from the reducers,
ordinary sample variations give full equation and cut transversality.
The residual $H$-orbit of a framed irreducible has positive dimension,
whereas its regular framed cut-down space would have dimension zero.
Thus no such irreducible point can remain.  Equivalently, the
irreducible quotient by $H$ has negative expected dimension and the
same local variations make its universal problem transverse.
There is no compactness boundary in this argument: a particle costs
an integer charge; a nonconstant lens-cylinder loss at total charge
$1/4$ leaves a flat cap, which misses the sphere cut.  These facts
follow from Lemma~\ref{lem:indices:lens-charge} and the compactness
of Lemma~\ref{lem:analysis:compactness}.  After the regular neighborhoods
of the reducer crossings are fixed, the remaining candidate set is
compact, so finitely many ordinary perturbation neighborhoods suffice.
The oriented sum of all points is therefore the winding sum.

There are only finitely many relevant main limits and finitely many
local branches there.  Cap adjustments can be supported in their
tangential parameter neighborhoods, away from the other true faces,
and copied to the flipped cap by tensoring.  They can also avoid
the legitimate instanton charts.  The established forbidden-stratum
exclusions persist under sufficiently small such adjustments.  To
see the required openness, a sequence of newly appearing forbidden
zeros would have an ideal limit by Lemma~\ref{lem:analysis:compactness}.
Further breaks and losses remain forbidden.  A fixed-type tuple has
a fixed-type limit; an unbroken reducible or active-line configuration
cannot converge smoothly to a legitimate irreducible zero-spinor
instanton because its parallel reduction persists.  These alternatives
are exactly the closed degenerations treated in
Proposition~\ref{prop:analysis:finite-induction}.  This proves the
compatibility asserted in the proposition.
\end{proof}

\subsection{The augmented operators at a matching pair}

Fix a main solution $u_M$ and a cap crossing $u_N$ as above, in a
single local branch.  All remaining metric parameters, cut marks,
and cap evaluation variables are included in the constrained
operators.  The neck length itself is held fixed during this
linearization.  Fix the effective phase on the main component by
the argument of a nonvanishing local sample function of original
phase character two.  Such a function exists at a free main solution.
Its differential is an isomorphism on the effective phase direction.

Let $G$ denote infinitesimal determinant-one gauge action and augment
the equation-and-cut derivative by $G^*$.  Use sufficiently high
little-H\"older completions; the field target has one less derivative
than the field domain.  On an end with coordinate $s\geq0$, a positive
weight means that $e^{\delta s}$ times the field is in the unweighted
space.  The number $\delta>0$ is less than the absolute value of every
nonzero eigenvalue of the boundary operator.  At the trace flat the
only zero mode of the augmented connection operator is the temporal
one-form $\tau\,ds$.  The adjoint spatial first cohomology vanishes,
the off-diagonal flat has no parallel sections, and the boundary
Dirac is invertible.

\begin{lemma}[The scalar obstruction and the weight change]
\label{lem:gluing:blocks}
After fixing phase as above, the positive-weight main augmented
constrained operator $L_M^+$ is an isomorphism.  On the cap the
positive-weight augmented constrained operator $L_N^+$ has zero
kernel and has precisely the one-dimensional cokernel
\begin{equation}
 \ell(f)=\int_N\langle\tau,f_{\mathrm{gauge}}\rangle.
 \label{eq:gluing:scalar-cokernel}
\end{equation}
Here $f_{\mathrm{gauge}}$ is its $G^*$-component and $\tau$ is the
parallel diagonal gauge generator on the cap.  Changing the cap
field weight from $+\delta$ to $-\delta$ gives an isomorphism $L_N^-$.
\end{lemma}

\begin{proof}
On the main completion gauge transformations may tend to parallel
constants.  Analytically their infinitesimal domain includes cutoff
extensions of these constants, in addition to decaying sections.
This extension is essential.  The gauge Laplacian satisfies
\[
 \langle G^*G\xi,\xi\rangle
 =\|d_a\xi\|_2^2+\|\xi\Phi\|_2^2.
\]
The boundary term is zero for an extended constant with decaying
derivative.  A vector in its kernel is parallel and annihilates the
nonzero main spinor, and is consequently zero.  In the complementary
weight, a homogeneous cokernel representative can have a constant
or linear parallel asymptote.  Pairing with the cutoff parallel
constants in the extended gauge domain forces the linear
coefficient, namely its boundary flux, to vanish.  The remaining
bounded representative satisfies the same integrated identity
and is zero.  One can
also obtain this directly on the cylindrical constant mode by
solving the scalar second-order equation with a free limiting
constant; the nonconstant boundary modes have the spectral gap.
Thus $G^*G$, on the extended gauge domain, is invertible.

It follows that the gauge augmentation gives a genuine slice for
unframed fields.  The main cut-down problem, with phase fixed, is
regular and isolated.  Its equation-and-cut derivative on that slice
is therefore an isomorphism.  Combining it with the gauge inverse
proves the assertion for $L_M^+$.

On the reducer cap, write $A_N$ for the reference connection.
The parallel gauge $\tau$ acts trivially on its field.  For a
variation $v=(a,\psi,\nu)$, including its mark $\nu$, integration
by parts gives
\[
 \ell(L_N^+v)=\int_N\langle\tau,d_{A_N}^*a\rangle=0
\]
for decaying variations; the spinor is zero.  The scalar gauge
equation consequently has the compatibility condition
\eqref{eq:gluing:scalar-cokernel}.  There is no other one.
Indeed, the cap gauge Laplacian $Q=G^*G$ on decaying gauge
parameters has zero kernel by integration.  In its complementary
weight, an adjoint homogeneous solution has a constant or linear
parallel asymptote; the other modes have the boundary spectral
gap.  On this one-ended cap, integration against the global
parallel $\tau$ forces the linear flux to vanish.  Integrating
the resulting bounded solution against itself shows that it is
parallel.  Thus the adjoint kernel is exactly $\R\tau$, and
$\im Q$ consists exactly of scalar targets of integral zero.

Given equation-and-cut target $y$, framed regularity from
\eqref{eq:gluing:cap-cut} and the invertible Dirac supply a field
and mark variation $u$ realizing $y$.  For any scalar target
$f_{\mathrm{gauge}}$ with integral pairing zero against $\tau$,
solve
\[
 Q\xi=f_{\mathrm{gauge}}-G^*u
\]
in decaying cap gauges.  Adding $G\xi$ corrects the scalar row
and leaves the equation-and-cut target unchanged, by equivariance
at the solution.  Conversely a zero of the augmented operator is
a framed gauge direction by constrained regularity, and then
$Q\xi=0$ makes it zero.  This gives zero kernel and full range
in $\ker\ell$.  Thus
$\ind L_N^+=-1$ with the full scalar target used here.

Only the eigenvalue zero is crossed in changing the weight.  Its
multiplicity is one, so $\ind L_N^-=0$.  A vector in the kernel of
$L_N^-$ has an expansion
\[
 a=c\tau\,ds+a_0,\qquad a_0=O(e^{-\epsilon s}),\quad\epsilon>\delta,
\]
with the same decay for derivatives and with decaying spinor; the
remaining finite-dimensional variables are the cap mark variables.
This expansion follows by the spectral decomposition on the end,
or by integrating each mode of the cylindrical operator.  Pair its
gauge equation with $\tau$ on the truncation $N_R$.  Since $\tau$ is
parallel, Green's formula gives
\[
 0=\int_{N_R}\langle\tau,d_{A_N}^*a\rangle
   =-\int_{\partial N_R}\langle\tau,a(\nu)\rangle.
\]
Letting $R$ tend to infinity forces
$c\,\|\tau\|_{L^2(\partial N)}^2=0$.  The vector hence lies in the
positive-weight domain and is zero by the preceding kernel assertion.
Index zero now proves surjectivity and the asserted inverse.
\end{proof}

Main parameters can also change the cap equations.  This does not
make the separated operator diagonal automatically.  The next
observation records the needed change of variables.

\begin{lemma}[Eliminating shared parameter derivatives]
\label{lem:gluing:triangular}
Let $p$ denote the main parameter and auxiliary-variable directions
which also occur in the cap equations or cap cut.  Locally the
limiting representative $\gamma$, collar gauges, and link metric
can be kept fixed.  In these choices, the limiting coupled derivative
can be reduced to $L_M^+\oplus L_N^-$ by a triangular change of variables.
The field part of its cap correction to $p$ decays exponentially.
\end{lemma}

\begin{proof}
Write $B_Np$ for the derivative of the cap equation and cut in these
directions.  The augmentation is defined at the reference field
and does not vary with $p$; thus $(B_Np)_{\mathrm{gauge}}=0$.
Moreover these variations do not change the limiting flat, so their
other components decay.  By \eqref{eq:gluing:scalar-cokernel},
$\ell(B_Np)=0$.  Lemma~\ref{lem:gluing:blocks} therefore gives a unique
positive-weight correction
\begin{equation}
 Kp=-(L_N^+)^{-1}_{\ker\ell}B_Np.
 \label{eq:gluing:parameter-correction}
\end{equation}
In this formula the inverse means the inverse from $\ker\ell$ to
the cap domain, including its mark variable.
Substituting $v_N=w_N+Kp$ changes the cap row
$L_Nv_N+B_Np$ into $L_Nw_N$.  The main row is unchanged, giving
the triangular elimination.  Elliptic regularity and the boundary
gap give exponential decay of the field component of $Kp$, with
derivatives, at some rate $\beta>0$ independent of a subsequent
choice of smaller $\delta$.  This uses the fixed boundary spectral
gap and the compact support or prescribed rapid decay of $B_Np$.
The mark component is an ordinary finite-dimensional correction
on the cap core.  All these
statements are uniform on the finite set of main limits and local
branches under consideration.
\end{proof}

\subsection{A finite-length inverse and the nonlinear gluing map}

\begin{proposition}[Regular gluing and exhaustion]
\label{prop:gluing:collar}
For every matching pair $(u_M,u_N)$ above and every sufficiently
large finite lens-neck length $D$, there is exactly one nearby
unbroken zero in a fixed local branch, modulo determinant-one gauge
and effective phase.  These zeros depend continuously, and smoothly
in finite-order solution charts, on $D$.  Every sequence of unbroken
zeros converging to this pair belongs to this collar for all
sufficiently large lengths.  The construction preserves the finite
branch lists, their multiplicities, and their weights.
\end{proposition}

\begin{proof}
Use an oriented cylinder $[0,D]\times\partial N$, with coordinate
$t$ increasing from the main body to the cap body.  Cut off the
exponentially decaying tails of $u_M$ and $u_N$ near the middle and
splice through $\gamma$.  Perform the same extensions for all
equation perturbations and cut representatives.  Lemma~\ref{lem:analysis:neck}
and the prescribed perturbation tails give a preglued field $u_D^0$
with residual, including the cuts,
\begin{equation}
 \|\mathcal F_D(0)\|\le C e^{-bD}
 \label{eq:gluing:residual}
\end{equation}
for some $b>0$, after decreasing $b$ to absorb fixed polynomial
factors.  Here $\mathcal F_D$ is the augmented, phase-fixed local
equation in field and finite-dimensional variables centered at
$u_D^0$.  The preglued fields have uniform bounds of every fixed
local differentiability order.  The perturbation library was
chosen with summable bounds at each such order.

\emph{Linear estimate.}
Choose $\delta>0$ below the boundary gap and as small as needed
below.  On the neck use the weight $w_D(t)=e^{\delta t}$, smoothly
continued as $1$ on the main core and $e^{\delta D}$ on the cap
core.  The cap's own mark variables and cut targets are scaled by
$e^{\delta D}$ as well.  The main variables are not scaled.
Viewed from the cap with outward coordinate $s=D-t$, the relative
weight is $e^{-\delta s}$; viewed from the main side it is
$e^{\delta t}$.  The limiting operators in these norms are
therefore exactly $L_M^+$ and $L_N^-$.
Figure~\ref{fig:gluing:weight} shows the two outward coordinates and
their induced weights.

\begin{figure}[htbp]
\centering
\begin{tikzpicture}[font=\small,>=Stealth]
 \draw[rounded corners=2pt,fill=blue!5] (0,0) rectangle (2,1.05);
 \draw[rounded corners=2pt,fill=blue!5] (8,0) rectangle (10,1.05);
 \draw (2,.2)--(8,.2) (2,.85)--(8,.85);
 \node at (1,.55) {main core};
 \node at (9,.55) {cap core};
 \node at (5,.55) {$w_D(t)=e^{\delta t}$};
 \draw[->] (2.2,1.4)--(7.8,1.4) node[midway,above] {$t$};
 \node[above] at (2,1.4) {$0$};
 \node[above] at (8,1.4) {$D$};
 \draw[->] (7.8,-.35)--(2.2,-.35)
    node[midway,below] {$s=D-t$};
 \node[align=center,below] at (1,-.15)
    {$w_D=1$\\$L_M^+$};
 \node[align=center,below] at (9,-.15)
    {$w_D=e^{\delta D}$\\$L_N^-$};
 \node[below] at (5,-1.05)
    {cap-relative weight: $e^{-\delta D}w_D=e^{-\delta s}$};
\end{tikzpicture}
\caption{The asymmetric weight on a finite lens neck.  Measured from
 the main core it induces the positive completion weight; after dividing
 by the cap-core factor $e^{\delta D}$ it induces the negative completion
 weight in the cap's outward coordinate.  The cap mark and cut target
 carry the same cap-core factor.}
\label{fig:gluing:weight}
\end{figure}

First transplant the triangular change of
Lemma~\ref{lem:gluing:triangular}, cutting off $Kp$ on the long
cylinder.  Its cutoff error is exponentially small, even with
the displayed weight, when $\delta$ is small compared with its
decay rate.  Write $\widetilde L_D$ for the resulting derivative.
There is a constant $C_\delta$, independent of large $D$, such that
\begin{equation}
 \|v\|_{X_{D,\delta}}
 \le C_\delta\|\widetilde L_Dv\|_{Y_{D,\delta}}.
 \label{eq:gluing:weighted-inverse}
\end{equation}
Here one may use the supremum of weighted local H\"older norms on
unit strips, with the usual norms on the compact cores and on the
finite-dimensional variables.  These are Banach norms on the
finite cylinder and give the same elliptic Fredholm problem.

We give an estimate proof to account for possible concentration in
the middle.  After conjugating by $w_D$, the model there is
$\partial_t+B-\delta$, with $B$ the self-adjoint augmented boundary
operator.  For every spectral value $\lambda$ its scalar inverse
on the full line integrates toward $+\infty$ when
$\lambda-\delta<0$ and toward $-\infty$ when
$\lambda-\delta>0$.  The respective kernels decay as
$e^{-|\lambda-\delta||t-s|}$.
Since $\dist(\delta,\operatorname{spec}B)>0$, spectral summation
and local elliptic estimates give a bounded full-cylinder inverse,
also on the uniformly local norms just specified.  In particular,
this estimate includes the parallel scalar mode $\lambda=0$.

Cover the glued manifold by the main completion model, the cap
completion model, and this cylindrical model.  Take cutoff
functions with transitions of a fixed, large width $R$, constant
on the compact cores and cut supports.  Apply the two inverse
estimates of Lemma~\ref{lem:gluing:blocks} and the cylinder estimate
to these cutoffs of $v$.  The commutators are at most
$C_\delta R^{-1}\|v\|$, in the corresponding one-order-lower norm.
Coefficient errors at the far ends tend to zero as $R$ increases;
errors from the actual splicing tend to zero as $D$ increases.
The compact graph terms have no unaccounted interaction between
the patches: terms supported on a body have its cutoff equal to
one, and tails or terms reaching the deep cylinder have arbitrarily
small exponentially weighted operator bounds.  The same tail
bounds control the finite-dimensional parameter forcing outside
its body's cutoff.  Shared parameter
terms have already been removed by
\eqref{eq:gluing:parameter-correction}.  Choose $R$ to absorb the
commutators and end errors, and then choose $D$ large to absorb the
remaining errors.  This proves
\eqref{eq:gluing:weighted-inverse}.  It also explains why merely
discarding the zero boundary eigenvalue would give an incorrect
estimate.

The unbroken constrained family has index one after fixing effective
phase.  Fixing the lens-neck length $D$ removes one parameter, so the
finite-length constrained problem has index zero.  Thus
\eqref{eq:gluing:weighted-inverse} implies that its derivative is
invertible.  Undoing the weight and the triangular change gives
an inverse $Q_D$ in the original norms with
\begin{equation}
 \|Q_D\|\le C e^{a\delta D}
 \label{eq:gluing:inverse-loss}
\end{equation}
for a fixed constant $a$.  A fixed polynomial in $D$ can be
absorbed by increasing $a$, for every fixed $\delta>0$.
No bound independent of $D$ in the original symmetric norms is
being asserted.

\emph{Nonlinear estimate.}
Set
$L_D=D\mathcal F_D(0)$ and
$R_D(v)=\mathcal F_D(v)-\mathcal F_D(0)-L_Dv$.
In a small ball of radius $r$ in the original field and variable
norms,
\begin{equation}
 \|R_D(v)-R_D(v')\|
 \le C e^{a'\delta D}r\,\|v-v'\|.
 \label{eq:gluing:nonlinear}
\end{equation}
The local quadratic field terms have this estimate by multiplication
in the chosen high-order spaces.  We spell out its moving-path
point because arbitrary same-order smoothness of translation is
not needed.  If $T_p$ denotes pullback along a smooth family of
paths or local transports, then with one derivative lost,
\[
 \|(T_p-T_0)h\|_{k-1,\alpha}
 \le C|p|\,\|h\|_{k,\alpha}.
\]
For the parameter derivative, subtract its value at the smooth
base field $u_D^0$ by first replacing the field by $u_D^0$ and
then replacing $p$ by zero.  The first difference is bounded by
$C\|u-u_D^0\|_{k,\alpha}$; the second is bounded by
$C|p|\|u_D^0\|_{k+1,\alpha}$.
Only the base field needs the extra derivative.  Differentiating
the transport ordinary differential equation gives these same
bounds for holonomies.  Products, smooth functions of finitely
many samples, localized outputs, and the energy-cutoff scalars
preserve them.  Consequently the derivative of the entire graph
formula differs from its derivative at the smooth preglued base
by at most $Cr$.  Integration on the segment from $v$ to $v'$
gives \eqref{eq:gluing:nonlinear}.  Little-H\"older completions
justify this first-derivative calculation; the compatible
higher-order solution charts are those of
Lemma~\ref{lem:analysis:scale}.  The prescribed summable
library bounds and the rescaling introduce at most the exponential
factor written in \eqref{eq:gluing:nonlinear}.

Choose $\delta$ so small that
$(2a+a')\delta<b/2$.  Put
$r_D=2C^2e^{-(b-a\delta)D}$, enlarging its fixed coefficient if
necessary.  Equations~\eqref{eq:gluing:residual}--\eqref{eq:gluing:nonlinear}
show that
\[
 v\longmapsto-Q_D\bigl(\mathcal F_D(0)+R_D(v)\bigr)
\]
maps the radius-$r_D$ ball into itself and has Lipschitz constant
less than $1/2$ for all large $D$.  It therefore has a unique fixed
point.  The same estimates give uniqueness on any ball of radius
$e^{-cD}$ for which
$(a+a')\delta<c<b-a\delta$; the inverse-scaled pregluing residual
$Q_D\mathcal F_D(0)$ then fits in that ball.  Dependence on finite $D$ follows by the
parameterized implicit theorem, with elliptic regularity giving
the compatible smooth solution charts.

\emph{Exhaustion.}
Let exact unbroken zeros tend to the specified matching pair while
$D\to\infty$.  There is no loss in their necks.  By
Lemma~\ref{lem:analysis:neck}, gauges on their two halves can be
chosen in which the deviation from $\gamma$ and all its derivatives
decay exponentially from the two bodies.  Flattening each half
near the middle then gives fields on the completed main and cap
pieces with exponentially small residuals.  Initially they are
$o(1)$-close to $u_M$ and $u_N$ in small positive decay weights:
on any fixed core this is convergence to the limit, and the
uniform exponential tail estimate controls the complementary
ends in any smaller weight.

Put these flattened fields in the fixed completed gauge slices
before applying the augmented estimates.  The completed main
gauge inverse and phase-fixing function give its slice; on the
cap the scalar gauge error has integral zero and the decaying
gauge inverse onto that range gives the framed slice.  These
are fixed-piece slice maps with constants independent of $D$.
Their gauge and phase rows now vanish exactly.  Equivariance
preserves the exponentially small equation-and-cut residuals.

The inverse estimate for $L_M^+$, followed by absorption of the
local nonlinear remainder, improves the main field and its
parameter differences from $o(1)$ to $O(e^{-b_1D})$ for some
$b_1>0$.  On the cap, $L_N^+$ is injective with closed range, so
it has the corresponding estimate on its domain even for a
residual not known beforehand to satisfy $\ell=0$.  Apply this
estimate to the cap difference after subtracting $Kp$; the already
controlled parameter difference and the flattening error are
exponentially small.  A second absorption therefore gives the
same exponential improvement on the cap.  Shrink $b_1$ if needed
to include all cutoff errors.

The half gauges may differ by an element of $H$.  The cap reducer
is fixed by $H$, so that mismatch is absorbed on the cap without
changing the matching point.  Splicing the half gauges makes the
global difference exponentially small.  Passing to the global
gauge slice costs at most a polynomial in $D$.  Here is a direct
bound: on a fixed main-core ball the limiting spinor has positive
norm bounded away from zero, and the gauge-action norm there
controls the norm of a Lie-algebra element.  Kato's inequality
and the one-dimensional Poincar\'e inequality along the tube give
\[
 \|\xi\|_2^2\le C(1+D)^2
 \bigl(\|d_a\xi\|_2^2+\|\xi\Phi\|_2^2\bigr).
\]
Local elliptic estimates upgrade this bound for the inverse of
$G^*G$ to the required field orders.  The small nonlinear slice
correction has the same polynomial loss, which exponential
closeness absorbs.  Hence the exact zeros lie in an
$O(e^{-b_2D})$-neighborhood of the pregluing for some $b_2>0$.
The fixed-piece half estimates give this positive rate before
the final small weight is selected.  Choose $\delta$ small
enough and then choose $c$ so that
\[
 (a+a')\delta<c<\min\{b-a\delta,b_2\}.
\]
The uniqueness ball of radius $e^{-cD}$ contains both the
constructed solution and all these exact zeros for large $D$.
They must therefore be the same zeros.  This proves
exhaustion, not just existence of a gluing map.

\emph{Branches and matching multiplicity.}
The cap crossing is fixed by $H$.  If both pieces are framed,
the diagonal $H$ quotient identifies every matching rotation
with the same pair.  If the main piece is unframed from the
outset, this is exactly the same quotient.  There is one glued
zero per cap crossing; no $H$-volume, eigenline choice, or extra
factor of two is introduced.  The effective phase was fixed on
the main piece, and the zero-spinor cap adds no phase parameter.

Each finite multisection is, in a local sample slice, a finite
list with multiplicities and weights.  At the separated pair
the active graph data belong to the appropriate parts.  A graph
which would compare frames across the true cut is absent there;
its finite-length tail is either zero or exponentially small.
Thus on each local branch the limiting equations are precisely
the matching branch equations used above.  The contraction
construction applies to each member of this finite list, with
the same thresholds after taking a maximum.  It commutes with
permutations of local lifts, and repeated branches keep their
multiplicities.  Accordingly it pulls the matching list back to
the list of unbroken zeros with exactly the original weights.
No new choice of branches, or averaging over matching rotations,
is made during gluing.
\end{proof}

\subsection{Completion of the compatible regularization}
\label{subsec:gluing:completion}

The lens collars constructed above complete the local description of
all permitted ends.  We can now regularize the remaining free stratum
relative to those collars and the instanton collars.

\begin{proof}[Proof of Theorem~\ref{thm:analysis:regularization}]
The ordinary construction is Proposition~\ref{prop:analysis:ordinary-library},
and its exact support and lower-level behavior follow from
Lemmas~\ref{lem:analysis:compactness} and \ref{lem:analysis:incidence}.
For ordinary non-PSC Floer ends the compactness and gradient-regularity
arguments are those in Section~\ref{sec:ordinary}; the isolated true-neck
argument is needed only for the present spinor family.
Proposition~\ref{prop:analysis:first-stage} verifies the fixed-type
requirements using single-valued data.  Lemma~\ref{lem:analysis:projection}
and Proposition~\ref{prop:analysis:finite-induction} exclude the forbidden
mixed candidates and give regular zero-dimensional trace-main problems.
The estimates and ideal limits are Lemmas~\ref{lem:analysis:compactness}
and \ref{lem:analysis:neck}.  Proposition~\ref{prop:analysis:instanton-link}
gives the legitimate instanton collars and their contribution.

Propositions~\ref{prop:gluing:regular-cap} and \ref{prop:gluing:collar}
now give exhaustive regular collars at the remaining no-particle lens
ends.  On the unbroken part use a countable small single-valued library
supported off these collars and the protected instanton collars, or
regularize compact remainders successively.  Universal surjectivity there
again uses its single $P$ component, and the expected dimension is one.
Such small final choices preserve the forbidden-stratum exclusions by
the compactness argument of
Proposition~\ref{prop:analysis:finite-induction}.

There are no other ideal boundaries by the numerical exclusions, now
with their analytic hypotheses proved.  The lens construction uses the
local cap operator and the exponential estimates, with no gluing at an
excluded positive-particle stratum.  With the instanton and lens collars
added, sequential compactness gives the claimed compact weighted
one-dimensional object.  Lemma~\ref{lem:analysis:stokes} supplies its
boundary rule.  All choices were made in arbitrarily small prescribed
norms and obeyed the support, parameter-locality, and tensor-identification
conditions from the start.
\end{proof}

\subsection{The local comparison of determinant orientations}

Table~\ref{tab:gluing:orientations} collects the conventions used in the
comparison.  The sign $s$ is the sign of the full real cap determinant
line, including its framing, mark, and cut factors.
\begin{table}[htbp]
\centering
\footnotesize
\begin{tabular}{@{}>{\raggedright\arraybackslash}p{.15\textwidth}>{\raggedright\arraybackslash}p{.24\textwidth}>{\raggedright\arraybackslash}p{.32\textwidth}>{\raggedright\arraybackslash}p{.19\textwidth}@{}}
\toprule
Datum & Ordinary convention & Spinor continuation & Reference \\
\midrule
Determinant and gauges & Determinant connection and transition maps fixed;
determinant-one gauges & Same fixed exterior data under a local cap
flip & Convention~\ref{conv:ordinary:index} \\
Framing & Unframed main; framed one-ended cap & The cap stabilizer
absorbs matching rotations, with no extra multiplicity &
Proposition~\ref{prop:gluing:collar} \\
Local dimensions & Interval: $1+1+1-3=0$ & Fixed framed cap:
$2+1-3=0$ & Theorem~\ref{thm:negative:unit};
Proposition~\ref{prop:gluing:regular-cap} \\
Real cap line & Ordered parameter, mark, cut, and framing factors;
comparison sign $s$ & The same full real comparison;
$\varepsilon_1=-s\varepsilon_0$ &
Proposition~\ref{prop:negative:local-sign} \\
Dirac factor & Absent from the ordinary operator & Canonical complex
orientation contributes the separate real sign $+1$ &
Proposition~\ref{prop:gluing:signs} \\
Phase and cuts & Ordinary instanton sign retained & Original character
two; effective phase quotients by $\{1,-1\}$; the link cut is
$\mathcal O(2)$ & Proposition~\ref{prop:analysis:instanton-link} \\
Branch weights & The integral count uses the fixed lift weights &
Lens matching preserves each branch weight; near instanton links the
total branch weight is one &
Propositions~\ref{prop:gluing:collar} and
\ref{prop:analysis:instanton-link} \\
\bottomrule
\end{tabular}
\caption{The ordinary conventions and their spinor continuation.
The interval count varies the metric and the sweep mark; the cap count
fixes the separating face and retains the framed cap directions.
The full comparisons, rather than a prescribed value of $s$, determine
the cancellation.}
\label{tab:gluing:orientations}
\end{table}

Let $\varepsilon(\mathbf e)$ denote the relative lift weights already
used in the integral ordinary composition, where
$\mathbf e=(e_1,\ldots,e_n)$.  They are fixed before constructing
the spinor problem.  The next proposition shows that these same
weights cancel the lens ends.

\begin{proposition}[The ordinary weights cancel the spinor lens ends]
\label{prop:gluing:signs}
For every regular trace-face main solution and every local branch,
the two caps obtained by flipping its lens bit have opposite
oriented contributions after multiplication by
$\varepsilon(\mathbf e)$.  This comparison is independent of the
main solution, respects all remaining gluings, and preserves the
branch weights.  In particular it is the relative sign of
Proposition~\ref{prop:negative:local-sign}, with no new sign choice.
\end{proposition}

\begin{proof}
Fix a regular trace-face main solution and one local branch.  The
exterior data agree under the lens-bit flip by
Proposition~\ref{prop:analysis:finite-induction}, and
Proposition~\ref{prop:gluing:regular-cap} chooses tensor-related cap
data.  We first recall the relative cap operation from
Proposition~\ref{prop:negative:local-sign} in these matching conventions.
If $x\in H^2(N;\Z)$ has $x(S)=1$, then
\[
 x^2=-\tfrac14,\qquad \PD(S)|_N=-4x.
\]
Tensoring on the cap by the line with first Chern class $-2x$
changes the rank-two determinant by $-4x$.
Globally, before this bit is flipped, $c\cdot S=0$ and
$S^2=-4$.  Thus
\begin{equation}
 \Delta c^2=-4,
 \qquad \Delta c_2=-1,
 \qquad \Delta\bigl(c_2-c^2/4\bigr)=0.
 \label{eq:gluing:flip-charge}
\end{equation}
The same flip adds $\PD(S)$ to $\Lambda$.  This is a local cap
replacement with the identical exterior bundle and identical
relative exterior charge; it is not a global choice of a
square root on the exterior.

The tensor line has the order-two flat boundary character.  In
the fixed common exterior determinant normalization, it
sends $\diag(i,-i)$ to $\diag(-i,i)$, a conjugate trace
representative.  Choose once a determinant-preserving parallel
boundary conjugation returning to the fixed $\gamma$.
Its ambiguity lies in the connected group $H$ and is absorbed
by the cap stabilizer.  At charge $1/4$ the relative cap degree
is fixed, and \eqref{eq:gluing:flip-charge} identifies the desired
cap sectors.  Use the intrinsic cap root obtained by tensoring
the old root with the twisting line, continued from the collapsed
sweep loop as in Proposition~\ref{prop:negative:local-sign}.
With these roots the normalized sphere path is unchanged before
boundary conjugation, which preserves its central target $-1$.
The adjoint data and sphere cut therefore correspond; copying
the cap curvature and cut perturbations gives the paired crossings
and full framed deformation-and-cut complexes of
Proposition~\ref{prop:gluing:regular-cap}.

Fix the ordering of parameter, mark, cut, and phase factors used
in the ordinary and spinor counts.  At a product-flat lens
splitting the instanton orientation line factors by determinant
gluing as
\begin{equation}
 \lambda_I(\mathrm{total})
 \cong \lambda_I(M;\mathrm{unframed})
       \otimes\lambda_I(N;\mathrm{framed}).
 \label{eq:gluing:excision}
\end{equation}
The framed cap factor includes the boundary-stabilizer orientation
which would otherwise occur in matching.  Use the full cap lines
$K_e$ and gluing isomorphisms
$\Gamma_e:\lambda_V\otimes K_e\longrightarrow\lambda_e$
of Proposition~\ref{prop:negative:local-sign}; $K_e$ includes the
framing, sweep-mark, and holonomy-target factors.  The comparison
and weights in \eqref{eq:negative:orientation-comparison} and
\eqref{eq:negative:lift-weights} are
\[
 \Phi=\Gamma_1(\id\otimes\tau_{\mathrm{cap}})\Gamma_0^{-1},
 \qquad \Phi(o_0)=s\,o_1,
 \qquad \varepsilon_1=-s\,\varepsilon_0,
\]
where $o_e$ are the previously chosen ordinary sector orientations.
Thus $\Phi(\varepsilon_0o_0)=-\varepsilon_1o_1$.
The map $\tau_{\mathrm{cap}}$ includes every framing and cut factor,
including the possible complex conjugation of an ordered
off-diagonal coordinate and reversal of the toral generator under
the boundary Weyl conjugation.  Only this full real determinant-line
sign $s$ is used.  The products of these comparisons are exactly
the weights $\varepsilon(\mathbf e)$ in the ordinary composition.

Proposition~\ref{prop:negative:local-sign} proves that this comparison
extends with constant sign over the relative configuration sectors:
the connection spaces are connected and the integral-lift orientation
systems have trivial monodromy on the determinant-one gauge quotient.
The fixed charge $1/4$ excludes a change of internal boundary degree.
The same proposition proves naturality under other gluings and
multiplicativity on disjoint caps; the constrained cap factors have
even parity.  These are comparisons of orientation torsors and require
no path of solutions.  They apply here because the exterior bundle,
relative degree, cap framing, and mark and cut ordering are precisely
the ones just identified.

Finally, homotope the spinor deformation operator, keeping its
symbol, to the direct sum of the instanton deformation operator
and the coupled complex Dirac operator.  Its orientation is the
instanton orientation times the complex Dirac orientation, followed
by the chosen parameter, cut, and effective-phase factors.
Tensoring the cap and applying the chosen boundary conjugation act
complex linearly on the spinor bundles, so the canonically oriented
Dirac determinant contributes the separate real sign $+1$.
This remains true if the interior Dirac operators are joined by a
homotopy rather than identified by an identical formula.  The exterior
phase, parameter, and all other cut factors agree in the paired
problems.  Thus the full spinor comparison has exactly the ordinary
relative sign above.

The perturbation prescription ties each exterior finite branch list
under the flip, rather than only its averaged value.  The chosen main
branch therefore has the same paired branch and rational weight.
Proposition~\ref{prop:gluing:collar} preserves those lists and weights
under gluing, with one match per cap crossing and no additional
$H$-multiplicity.  The opposite orientation comparison consequently
holds for each weighted lens-boundary contribution.
\end{proof}

\subsection{Closing the one-dimensional count}

\begin{theorem}[Contradiction from the boundary count]
\label{thm:gluing:contradiction}
The finite stack of Theorem~\ref{thm:stack:nonzero} cannot exist
under an orientation-preserving identification $Y_{-2}\cong Y_2$.
Consequently the remaining oriented cosmetic pair of slopes
$-2$ and $2$ is impossible.
\end{theorem}

\begin{proof}
Let $\Omega\ne0$ be the oriented ordinary instanton count of
Theorem~\ref{thm:stack:nonzero}, with the fixed lift weights
$\varepsilon(\mathbf e)$ and the prescribed cap insertions.
Put
\[
 r=n_D^{\mathrm{tot}}>0,\qquad \eta=r-1.
\]
The ordinary dimension equation is
$D_I^{\mathrm{tot}}+n-2n-z=0$.  Adding the spinor directions,
dividing by effective phase, and imposing the $\eta$ complex
phase cuts gives
\begin{equation}
 D_I^{\mathrm{tot}}+2r+n-2n-z-1-2\eta=1.
 \label{eq:gluing:one-dimensional}
\end{equation}
By Theorem~\ref{thm:analysis:regularization} the free zero set
is therefore an oriented one-dimensional finite weighted
branch space, with the exclusion and compactness properties
already proved there.

The instanton-link contribution is supplied by
Proposition~\ref{prop:analysis:instanton-link}.  At a legitimate
regular cut-down instanton the phase quotient of the normal link is
$\CP^{r-1}$, and the $\eta=r-1$ original-phase-character-two cuts
are sections of $\mathcal O(2)$.  With $h=c_1(\mathcal O(1))$,
\begin{equation}
 \left\langle(2h)^{r-1},[\CP^{r-1}]\right\rangle
 =2^{r-1}=2^\eta.
 \label{eq:gluing:link-multiplicity}
\end{equation}
For $r=1$ this means the single point of $\CP^0$ and no phase cuts.
The proposition gives the complex orientation multiplied by the ordinary
instanton sign, up to one common boundary-orientation convention.
The mixed and final interior perturbations vanish in these charts,
and their repeated zero branches have total weight one.
Consequently the total instanton-link boundary contribution is
$\sigma_0 2^\eta\Omega$, for one common
$\sigma_0\in\{1,-1\}$.

Choose small disjoint radial neighborhoods of these finitely
many instantons, and cut each of the finitely many regular lens
collars at a sufficiently large finite $D$.
Proposition~\ref{prop:gluing:collar} identifies every end of the
latter type with its matching main point and cap crossing.
The remaining part of the free zero set is compact.  Indeed,
an escaping sequence has an ideal limit by
Lemma~\ref{lem:analysis:compactness};
Theorem~\ref{thm:indices:exclusion} and
Proposition~\ref{prop:analysis:finite-induction} exclude every
limit except the neighborhoods just removed.  The exhaustion
assertion of Proposition~\ref{prop:gluing:collar} and the instanton
blow-up charts preclude escape even to those limits from the
truncated remainder.

Apply Lemma~\ref{lem:analysis:stokes} to this compact oriented
one-dimensional remainder.  Its finite local branch lists have the
compatible weighted germs proved there, and the instanton and lens
collars preserve their weights.  Branches extend by the same repeated
zero list across a perturbation-support boundary, so no artificial
boundary point is introduced.  The total weighted oriented boundary
is therefore zero.

Every lens contribution cancels with its bit-flipped partner
by Proposition~\ref{prop:gluing:signs}, using the same
$\varepsilon(\mathbf e)$ as in $\Omega$.  The exclusions have
left no other boundary.  Stokes therefore gives
\[
 0=\sigma_0 2^\eta\Omega.
\]
This contradicts $\Omega\ne0$ over $\Q$.  All neighborhoods
used here are the regular zero-particle neighborhoods proved
above; no gluing assertion at a positive-particle level has
entered the boundary argument.
\end{proof}

\begin{proof}[Proof of Theorem~\ref{thm:main}]
The reduction in Section~\ref{sec:introduction} shows that a
distinct orientation-preserving cosmetic pair on a nontrivial
knot in $S^3$ would have slopes $-2$ and $2$ and genus two.
The meridional slope is already excluded by the absence of a
second sphere filling, and zero cannot pair with a different
slope because of first homology.  An oriented homeomorphism
in the remaining case can be taken to be a diffeomorphism,
giving the identification used throughout the construction.
Theorem~\ref{thm:gluing:contradiction} excludes that case.
Thus any orientation-preserving homeomorphism between two
fillings in the statement has equal slopes.
\end{proof}
